% !TEX root = ../more_for_less.tex

The genomic revolution has fundamentally changed our understanding of human history \citep{green2010draft, gravel2011demographic, tennessen2012evolution, lazaridis2014ancient, haak2015massive}, the genetics of complex disease \citep{lek2016analysis, karczewski2020mutational, wang2021rare, backman2021exome,chen2023genomic}, and the development of cancer \citep{cancer2008comprehensive, cancer2013cancer, kandoth2013mutational, jayasinghe2018systematic, hoadley2018cell, huang2018pathogenic}. As more and more sequencing data is gathered, novel variants continue to be discovered \citep{lek2016analysis, karczewski2020mutational} and yield a deeper insight into the evolutionary processes that have shaped the observed patterns of genetic variation \citep{agarwal2021mutation, agarwal2023relating}. Moreover, rare variants, which are typically detected in large samples, have become increasingly important in understanding the etiology of disease and are a key tool in development of novel therapeutics against common disease \citep{szustakowski2021advancing}.

Most studies of genetic variation sample from two or more groups that may share some portion of their variation, such as different human populations or different tumor types. For instance, sampling from multiple human populations is essential for discovering more variants \citep{karczewski2020mutational,chen2023genomic} and is crucial to expanding equity and inclusion in genomic medicine \citep{popejoy2016genomics, hindorff2018prioritizing, fatumo2022roadmap}. Due to the shared evolutionary history of all humans, in this setting many common variants will be shared between populations \citep{biddanda2020variant}, but differences in allele frequency between populations are key to understanding adaptation \citep{novembre2009spatial} and disease \citep{fu2011multi, assimes2021large, roselli2018multi}. Similarly, when collecting disease case and control samples to examine the enrichment of rare variants in cases, the vast majority of the variants will be shared. Rare variants shared across different type of diseases have drawn attention since they may be linked to common genetic bases of different diseases \citep{rafnar2009sequence, bojesen2013multiple, rashkin2020pan}. On the other hand, most variation between tumors in different individuals will be private, but shared variation between tumors across samples may point to the genetic basis of cancer, and thus paths toward novel treatments \citep{st2014genomic, ma2015rise,prior2020frequency,pirozzi2021implications}.


Despite large decreases in the cost of sequencing \citep{christensen2015assessing, park2016trends, schwarze2018whole}, detecting novel variants still requires a large enough set of samples to benefit from careful planning. This is especially true when attempting to expand our catalogue of variation across diverse human groups, because some of the most diverse human groups have been studied the least \citep{popejoy2016genomics}. It would therefore be of significant practical utility to have a reliable and data-informed planning method; namely, a prediction method should accurately estimate the number of novel variants in a new (follow-up) study of multiple populations, given data from a typically small, previous (pilot) study. Moreover, prediction of the number of shared or private variants seen in a follow-up study can be useful as a null model for identifying ``unusual'' variants or classes of variants that may indicate their role in biologically interesting processes, such as adaptation or disease. In all of these cases, it is crucial to account for novel variants that are shared between groups as opposed to those that are private to a single group.

Several researchers have developed approaches to predict the number of new variants that will be detected in a follow-up sequencing study, given information from a pilot study \citep{Camerlenghi2021, Masoero2022, chakraborty2019using, zou2016quantifying,gravel2014predicting,ionita2009estimating}.  However, all of these methods assume samples are from a single population. So these methods cannot be used directly to predict the number of new shared variants across different populations in a follow-up study.
And because these methods do not model the possibility of shared variants between groups, single-population approaches can be expected to misestimate the number of novel variants in a follow-up study.  A complementary body of work models variant discovery within and between groups as a way to guide rational imputation panel design \citep{huang2013genotype, huang2009genotype, kang2015choosing, jewett2012coalescent, zhang2013genotype}. Although these models consider variant sharing between groups, they are not designed to predict new variant discovery, but rather to examine the ability to impute already-known variants in new samples; nonetheless, it reveals that it is important to quantify both the amount of variation that is shared between groups as well as the variation that is private to any given group. 


In this work, we develop an approach to predict the number of novel variants discovered in samples from multiple populations or groups. We focus on the case of two populations here, although in principle our approach can be extended to more than two. We work in the setting where a pilot study has collected genetic data in each of the two groups, and the experimenter wishes to collect additional data in a larger study.
The variant prediction challenge can be tackled in two ways: (1) one could construct an explicit model of variation within and between groups based on biological principles or (2) one could employ a ``black box'' that is agnostic to the underlying biology. This work, like past single-population approaches, adopts the latter approach. We hope that, by taking the black-box approach, our method can apply to data that arise from very different processes -- ranging across tumor sequencing, case-control studies, population genetics, and other settings. 

In the one-population case, \citet{Masoero2022} showed that a Bayesian nonparametric (BNP) approach can deliver strong empirical performance. 
However, we here prove that a natural two-population extension to the method of \citet{Masoero2022} fails for fundamental reasons (\cref{sec:cant_have_it_all}). We nonetheless demonstrate how to take advantage of the flexibility of the Bayesian framework (\cref{sec:bnp_two_pops}) to develop a novel two-population BNP model (\cref{sec:our_model}) and methodology (\cref{sec:prediction_all}). We provide theory to establish the desirable power-law properties of our model in \cref{sec:powerlaw}. Finally, we use both simulations and analysis of real cancer and population genetic datasets to showcase the accuracy and robustness of our model (\cref{sec:experiments}); along the way, we highlight multi-population circumstances where using one-population approaches can be misleading.


\paragraph{Additional related work.}
As in the work of \citet{Camerlenghi2021, Masoero2022, chakraborty2019using, zou2016quantifying,gravel2014predicting,ionita2009estimating}, our work considers a data point as the genetic sequence of an individual, and our model allows that each data point can exhibit multiple variants. As observed by \citet{Masoero2022}, this setup can be view as an example of a ``feature allocation'' problem \citep{ghahramani2005infinite,Broderick2013} in that each data point (an individual's genetic sequence) can belong to zero, one, or more groups, a.k.a.\ features (where each feature corresponds to the presence of a particular variant). \minorrevision{Feature allocations are not specific to applications in genetics; they have been applied, for instance, in ecology \citep{ghilotti2024bayesian}.}
A closely related but distinct setting with a rich history in the Bayesian nonparametric literature is the partition (a.k.a.\ species sampling) setting, in which each data point must belong to exactly one group, or partition element. The works of \citet{lijoi2007bayesian} and \citet{favaro2009bayesian} are important milestones in the Bayesian nonparametric literature, as they first proposed Bayesian nonparametric methods for a classical problem in ecology: predicting the number of new species in follow-up study given observations from a pilot study \citep{efron1976estimating}.
Other related work in the species sampling literature includes \citet{favaro2012new}, which developed a Bayesian nonparametric model to estimate the chance of observing a particular group in one or more follow-up samples of any given frequency in the pilot samples, and \citet{favaro2016rediscovery}, which provides a theoretical connection between the classic Good-Turing estimator \citep{good1953population} on discovery probability and the Bayesian nonparametric approaches of \citet{lijoi2007bayesian, favaro2009bayesian, favaro2012new}. 
Further, \citet{camerlenghi2017bayesian} propose predicting the number of new species to be observed in an enlarged sample using the hierarchical Pitman-Yor process, a hierarchical model that accounts for heterogeneity among individuals in species sampling cases. 
Most recently, \citet{camerlenghi2019distribution} provided a theoretical study for a broader class of hierarchical species sampling models. 
These works often find application in genetics, particularly in discovering genes in cDNA libraries.
In particular, in their setting, each sample is a cDNA read representing the expression of a single gene.
Thus, each read belongs to one and only one partition element (gene). 
This setup stands in contrast to the present paper, where each data point is a genetic sequence from an organism
and may exhibit an arbitrary non-negative integer number of features (genetic variants).

